% !PRINTED 281
\weberpage{281}
Es sei
\weq{1}{f(x) = a_0 x^n + a_1 x^{n-1} + \cdots + a_{n-1} x + a_n = 0}
die vorgelegte Gleichung und
\weq{2}{x_1,\, x_2 \ldots x_n}
ihre Wurzeln, die wir \gesperrt{von einander verschieden annehmen}.

Wir benutzen die in \S.\,68 definirten Functionen
\weq{3}{f_0(x),\, f_1(x),\, f_2(x) \ldots f_{n-1}(x)}
und setzen wie dort, indem wir unter $t_0, t_1 \ldots t_{n-1}$ unabh\"angige Variable verstehen,
\weq{4}{y = t_{n-1} f_0(x) + t_{n-2} f_1(x) + \cdots + t_1 f_{n-2}(x) + t_0 f_{n-1}(x).}
Es m\"ogen $y_1, y_2 \ldots y_n$ die Werthe sein, die $y$ f\"ur $x = x_1, x_2 \ldots x_n$ annimmt.

Ist nun $\alpha$ ein beliebiger reeller Werth, so setzen wir
\weq{5}{H_\alpha = (x_1-\alpha)y_1^2 + (x_2-\alpha)y_2^2 + \cdots + (x_n-\alpha)y_n^2.}
Dies ist eine quadratische Form der $n$ Variablen $t$, und wir wollen zun\"achst die Zahl ihrer negativen und positiven Glieder bestimmen. Ist $x_1$ eine reelle Wurzel, also auch $y_1$ reell, so ist das Quadrat $(x_1-\alpha)y_1^2$ positiv oder negativ, je nachdem $x_1$ gr\"osser oder kleiner als $\alpha$ ist. Bilden aber $x_1, x_2$ ein imagin\"ares Paar, so sind auch $y_1, y_2$ conjugirt imagin\"ar, und
\[ (x_1-\alpha)y_1^2 + (x_2-\alpha)y_2^2 \]
zerlegt sich in ein \gesperrt{positives} und ein \gesperrt{negatives} Quadrat. Dies ergiebt sich wie im \S.\,79, wenn man $y_1\sqrt{x_1-\alpha} = u+iv$, $y_2\sqrt{x_2-\alpha} = u-iv$ setzt.

Daraus folgt, dass die Anzahl $N_\alpha$ der negativen Quadrate in $H_\alpha$ gleich ist der Anzahl der imagin\"aren Paare, vermehrt um die Anzahl der reellen Wurzeln, die kleiner als $\alpha$ sind.

Nehmen wir also eine zweite reelle Zahl $\beta > \alpha$, bilden die Function $H_\beta$ und bezeichnen mit $N_\beta$ die Anzahl ihrer negativen Quadrate, so ist die \gesperrt{Differenz}
\[ N_\beta - N_\alpha \]
\gesperrt{gleich der Anzahl der reellen Wurzeln zwischen $\alpha$ und $\beta$.}

Man erh\"alt also ein Mittel zur Bestimmung dieser Zahl, d.\,h. zur L\"osung des Sturm'schen Problems, wenn man die Co\"efficienten der Function $H_\alpha$ als Functionen der Co\"efficienten von $f(x)$ und von $\alpha$ darstellt, und dann die Zahl $N_\alpha$ untersucht,
% [sic] The printed page ends the complete sentence before § 91 with a comma.
